%%%%%%%%%%%%%%%%%%%%%%%%%%%%%%%%%%%%%%%%%%%%%%%%%%%%%%%%%%%%%%%%%%%%%%%
%% OUTLINE / SKELETON -- Hội thảo quốc gia lần thứ XXIX (VNICT 2026)
%% Giới hạn: 6 trang. Mỗi \section bên dưới kèm ngân sách trang/chữ.
%% Preamble giữ nguyên theo bare_conf.tex của BTC -- KHÔNG sửa lề/font.
%%%%%%%%%%%%%%%%%%%%%%%%%%%%%%%%%%%%%%%%%%%%%%%%%%%%%%%%%%%%%%%%%%%%%%%
\documentclass[10pt, conference, a4paper, compsocconf]{IEEEtran}
\usepackage{amsmath,amsxtra,amssymb,amsthm,latexsym,amscd,amsfonts}
\usepackage[utf8]{vietnam}

\usepackage[top=2.7cm, bottom = 5.4cm, left=2.72cm, right = 3.0cm]{geometry}
\headsep = 0.7cm
\headheight = 1.0cm
\footskip = 1.0cm
\voffset = -0.74 cm

\usepackage{fancyhdr}
\pagestyle{fancy}
\renewcommand{\sectionmark}[1]{\markright{\MakeUppercase{#1}}{}}
\renewcommand{\headrulewidth}{0pt}

\ifCLASSINFOpdf
  \usepackage[pdftex]{graphicx}
  \DeclareGraphicsExtensions{.pdf,.jpeg,.png,.jpg}
\else
  \usepackage[dvips]{graphicx}
  \DeclareGraphicsExtensions{.eps}
\fi

\usepackage{array}
\usepackage[font=footnotesize]{subfig}
\usepackage{booktabs}
\usepackage{algorithmic}

\makeatletter
\def\ps@IEEEtitlepagestyle{%
\def\@oddhead{\centering{\mbox{\footnotesize{\textit{\;\;\;\; Hội thảo quốc gia lần thứ XXIX: Một số vấn đề chọn lọc của Công nghệ thông tin và truyền thông -- Hà Nội, 7-8/11/2026}}}}%
\def\@evenhead{\centering{\mbox{\footnotesize{\textit{\;\;\;\; Hội thảo quốc gia lần thứ XXIX: Một số vấn đề chọn lọc của Công nghệ thông tin và truyền thông -- Hà Nội, 7-8/11/2026}}}}}%
\def\@oddfoot{\scriptsize \thepage \hfil }%
\def\@evenfoot{\scriptsize \hfil \thepage}
}
}
\makeatother

\begin{document}
\pagenumbering{gobble}
\fontsize{10}{12}
\selectfont

\fancyhead[RE,LO]{\centering{\footnotesize{\textit{Hội thảo quốc gia lần thứ XXIX: Một số vấn đề chọn lọc của Công nghệ thông tin và truyền thông -- Hà Nội, 7-8/11/2026}}}}

%%%%%%%%%%%%%%%%%%%%%%%%%%%%%%%%%%%%%%%%%%%%%%%%%%%%%%%%%%%%%%%%%%%%%%%
% TITLE -- 3 phương án, chọn 1:
%  (a) Quy hoạch trạm sạc xe điện có xét ràng buộc lưới điện phân phối
%      bằng học tăng cường sâu trên đồ thị đô thị
%  (b) Học tăng cường sâu dựa trên mạng nơ-ron đồ thị cho bài toán mở
%      rộng mạng lưới trạm sạc xe điện tại Hà Nội
%  (c) Đặt trạm sạc xe điện nhận biết lưới điện: tiếp cận MDP và GNN-DQN
%%%%%%%%%%%%%%%%%%%%%%%%%%%%%%%%%%%%%%%%%%%%%%%%%%%%%%%%%%%%%%%%%%%%%%%
\title{Quy hoạch trạm sạc xe điện có xét ràng buộc lưới điện phân phối\\
bằng học tăng cường sâu trên đồ thị đô thị}

\author{\IEEEauthorblockN{Nguyễn Văn A\IEEEauthorrefmark{1}, Nguyễn Văn B\IEEEauthorrefmark{2}}
	\IEEEauthorblockA{\IEEEauthorrefmark{1}Đơn vị 1, \IEEEauthorrefmark{2}Đơn vị 2}
	\IEEEauthorblockA{Email: a@abc.xyz, b@abc.xyz}
}

\maketitle

%======================================================================
\begin{abstract}
% NGÂN SÁCH: 150--180 từ (~1/4 cột). Cấu trúc 6 câu:
% 1. Bối cảnh: xe điện tăng nhanh ở đô thị VN, mở rộng mạng trạm sạc là nút thắt.
% 2. Khoảng trống: các nghiên cứu đặt trạm thường chỉ tối ưu phủ/nhu cầu, bỏ qua
%    khả năng tải còn lại của lưới phân phối -> kế hoạch "tốt trên giấy" nhưng
%    gây quá tải thanh cái 22 kV.
% 3. Đề xuất: mô hình hoá bài toán thành MDP mở rộng kế hoạch hiện hữu, hàm
%    thưởng gồm lợi ích xã hội, chi phí đi lại/chờ/sạc, độ công bằng và phạt
%    ràng buộc lưới; chính sách học bằng DQN với bộ trích đặc trưng đồ thị.
% 4. Dữ liệu: 4 quận Hà Nội (Cầu Giấy, Đống Đa, Nam Từ Liêm, Tây Hồ), đồ thị
%    đường bộ OSM + lưới điện dựng từ toạ độ TBA 500/220/110/22 kV thực tế.
% 5. Kết quả: vượt greedy-benefit / greedy-demand / GA về điểm chuẩn hoá, đồng
%    thời đưa số thanh cái quá tải về 0 (điền số cụ thể X%, Y).
% 6. Ablation khẳng định vai trò của phạt lưới và mô hình nhu cầu suy giảm.
% LƯU Ý BTC: không dùng ký hiệu toán/ký tự đặc biệt trong tiêu đề và tóm tắt.
\end{abstract}

\begin{IEEEkeywords}
trạm sạc xe điện; học tăng cường sâu; mạng nơ-ron đồ thị; lưới điện phân phối;
bài toán định vị cơ sở; quy hoạch hạ tầng đô thị
\end{IEEEkeywords}

\IEEEpeerreviewmaketitle

%======================================================================
\section{Giới thiệu}
% NGÂN SÁCH: 0,75 trang (~1300 từ đôi cột => ~600--650 từ). 4 đoạn.
%
% Đ1 (bối cảnh, 5--6 câu): tốc độ điện hoá phương tiện; hạ tầng sạc công cộng
%   là điều kiện cần; đặc thù đô thị VN (mật độ dân cư cao, quỹ đất/giá đất là
%   ràng buộc mạnh, lưới phân phối đã vận hành gần ngưỡng).
%
% Đ2 (vấn đề, 4--5 câu): đặt trạm KHÔNG chỉ là bài toán phủ nhu cầu.
%   - Mỗi trạm là một phụ tải mới quy về thanh cái 22 kV gần nhất;
%   - Dung lượng dự phòng của thanh cái là hữu hạn và mang tính CỘNG DỒN:
%     nhiều trạm "tốt" cục bộ vẫn có thể cùng làm quá tải một thanh cái;
%   - Khoảng cách tới điểm đấu nối cũng là chi phí (chi phí đường dây/tổn thất).
%   Đây là loại ràng buộc mà các phương pháp tham lam/tĩnh xử lý kém vì nó phụ
%   thuộc toàn bộ chuỗi quyết định trước đó.
%
% Đ3 (định vị đóng góp, 3--4 câu): tách khỏi (i) nhóm định vị cơ sở cổ điển
%   (MILP/GIS, bỏ qua lưới) và (ii) nhóm phân tích hosting capacity (có lưới
%   nhưng tối ưu tĩnh, một lần). Ta xử lý bài toán như ra quyết định TUẦN TỰ có
%   ngân sách suy giảm, xuất phát từ kế hoạch trạm HIỆN HỮU chứ không từ số 0.
%
% Đ4 (đóng góp -- itemize 4 gạch đầu dòng, mỗi gạch 1--2 dòng):
%   1) Phát biểu MDP cho bài toán mở rộng mạng trạm sạc, hợp nhất mục tiêu phúc
%      lợi xã hội với ràng buộc dung lượng lưới phân phối ở mức thanh cái.
%   2) Mô hình nhu cầu động, suy giảm theo khoảng cách và theo công suất đã lắp,
%      khiến lợi ích biên của trạm mới bão hoà -> tránh đặt trùng lặp.
%   3) So sánh có kiểm soát 4 cách mã hoá trạng thái (MLP, MLP+đặc trưng đồ thị,
%      GCN, attention+PMA) trên cùng một thuật toán DQN và cùng một môi trường.
%   4) Đánh giá thực nghiệm trên 4 quận Hà Nội với dữ liệu TBA thực tế, đối
%      sánh 3 baseline và 4 kịch bản ablation.
%
% Đ5 (1 câu): bố cục bài báo.

%======================================================================
\section{Các nghiên cứu liên quan}
% NGÂN SÁCH: 0,5--0,6 trang (~450 từ). 3 đoạn, KHÔNG chia mục con (tiết kiệm chỗ).
%
% Đ1 -- Định vị trạm sạc: p-median/MCLP, MILP, nghiên cứu GIS đa tiêu chí.
%      Điểm chung: mục tiêu phủ/nhu cầu/chi phí, coi lưới điện là vô hạn.
% Đ2 -- Metaheuristic và học tăng cường cho quy hoạch hạ tầng: GA/PSO cho định
%      vị; DRL cho bài toán tổ hợp trên đồ thị (S2V-DQN và các biến thể), DRL
%      cho định giá/điều độ sạc. Chỉ ra: DRL cho SITING (không phải routing hay
%      pricing) còn ít, đặc biệt là dạng mở rộng kế hoạch hiện hữu.
% Đ3 -- Ràng buộc lưới: hosting capacity, siting có ràng buộc dòng công suất,
%      đánh giá tác động phụ tải sạc lên lưới hạ thế.
% Câu kết (2 câu): khoảng trống = ra quyết định tuần tự + ràng buộc lưới cộng dồn
%      + biểu diễn đồ thị đô thị, đồng thời trong một khung duy nhất.

%======================================================================
\section{Phát biểu bài toán}
\label{sec:problem}
% NGÂN SÁCH: 1,0--1,1 trang. Đây là phần "lõi khoa học", đừng cắt.

\subsection{Mô hình hệ thống}
% ~250 từ + 1 bảng ký hiệu nhỏ (tuỳ chọn, cắt nếu thiếu chỗ).
% - Đồ thị đường bộ G=(V,E) trích từ OSM cho từng quận; mỗi nút v mang:
%   toạ độ (x,y), dân số pop, nhu cầu cơ sở demand, giá đất land_price,
%   street_count, tỉ lệ trạm sạc tư nhân private_cs. Cạnh mang chiều dài.
% - Trạm sạc s = (vị trí, cấu hình bộ sạc, thuộc tính dẫn xuất). 12 mức công
%   suất bộ sạc {3,...,250} kW với suất đầu tư tương ứng; tối đa K=100 trụ/trạm.
% - Kế hoạch hiện hữu P_0 (các trạm đã có) là điểm xuất phát; ngân sách B.
% - Lưới điện: thanh cái 500/220/110/22 kV dựng từ toạ độ TBA thực tế của Hà
%   Nội; mỗi nút đồ thị ánh xạ tới thanh cái 22 kV gần nhất; mỗi thanh cái có
%   dung lượng khả dụng available_mw sau khi trừ phụ tải nền.

\subsection{Hàm mục tiêu phúc lợi xã hội}
% ~300 từ + 3 phương trình. ĐÂY LÀ PHẦN PHẢI VIẾT ĐÚNG CODE.
% (1) Nhu cầu động: d_v(P) = d_v^0 * exp(-eta * sum_{s: dist<r_s} c_s e^{-beta*dist})
%     với eta = 0.4 (scaling), beta = 0.6 (distance decay).  <-- ĐÓNG GÓP 2
% (2) Thời gian chờ: hàng đợi M/M/1/N tại mỗi trạm, lambda = tổng xe quy đổi từ
%     dân số (ev_per_capita = 0.022), mu = capability*1000/BATTERY,
%     N_hiệu_dụng = max(100, ceil(lambda)); W_s theo định luật Little.
% (3) Điểm chuẩn hoá:
%     cost = (alpha*travel + (1-alpha)*(charging + waiting))/3, alpha = 0.8
%     fairness = 1/(1 + std(số trạm phủ mỗi nút))
%     score = (benefit - cost + fairness)/3 - w_grid * penalty_grid
% (4) penalty_grid = phạt khoảng cách đấu nối (trung bình theo trạm) + phạt
%     vượt dung lượng CỘNG DỒN theo từng thanh cái.
\begin{equation}
\label{eq:score}
\mathcal{S}(P) = \frac{B(P) - C(P) + F(P)}{3} - w_g\,\Phi_{\text{grid}}(P)
\end{equation}

\subsection{Phát biểu MDP}
% ~350 từ. Định nghĩa hình thức:
% - Trạng thái s_t: ma trận đặc trưng nút (N x 10) gồm lon, lat, nhu cầu tĩnh,
%   nhu cầu động, giá đất, street_count, công suất trạm tại nút, khoảng cách
%   tới lưới, dung lượng còn lại của thanh cái cục bộ, khoảng cách tới trạm gần
%   nhất; cộng 3 đặc trưng toàn cục (ngân sách còn lại, tiến độ episode,
%   độ cải thiện điểm so với đầu episode). Tất cả co giãn về [-1,1].
% - Hành động a_t in {0,...,4}: MACRO-ACTION tham số hoá bằng heuristic --
%   (0) mở trạm mới tại nút lợi ích biên cao nhất, (1) mở trạm mới tại nút nhu
%   cầu chưa đáp ứng cao nhất, (2)/(3) bổ sung trụ sạc cho trạm hiện có theo
%   lợi ích / theo nhu cầu, (4) di dời một trụ từ trạm kém hiệu quả sang trạm
%   thiếu. NHẤN MẠNH: đây là lựa chọn thiết kế cốt lõi -- nó giữ không gian hành
%   động ở kích thước 5 thay vì O(|V| x 12), giúp bài toán khả học; heuristic
%   quyết định "ở đâu", chính sách quyết định "chiến lược nào, khi nào".
% - Chuyển trạng thái: trừ ngân sách, cập nhật gán nút-trạm, cập nhật phụ tải
%   thanh cái. Kết thúc khi hết ngân sách hoặc t >= |V|/2.
% - Hàm thưởng (viết đúng như evaluation() trong code):
%     r_t = [S(P_t) - S(P_{t-1})] + w_b * max(0, S(P_t) - S*_{t-1})
%   Giải thích 2 câu: số hạng 1 là potential-based shaping (tổng telescoping =
%   cải thiện ròng của episode); số hạng 2 thưởng riêng cho việc phá kỷ lục, do
%   ta chỉ quan tâm KẾ HOẠCH TỐT NHẤT tìm được chứ không phải trạng thái cuối.
\begin{equation}
\label{eq:reward}
r_t = \big[\mathcal{S}(P_t)-\mathcal{S}(P_{t-1})\big] + w_b\max\big(0,\;\mathcal{S}(P_t)-\mathcal{S}^{*}_{t-1}\big)
\end{equation}

%======================================================================
\section{Phương pháp đề xuất}
% NGÂN SÁCH: 1,2--1,3 trang, gồm 1 hình kiến trúc (Fig.~1) + 1 giả mã.

\subsection{Mã hoá quan sát trên đồ thị}
% ~200 từ. Vì sao cần chuẩn hoá theo hằng số riêng của từng quận
% (scaling_constants.json): các quận khác nhau về quy mô/dân số/giá đất, chuẩn
% hoá riêng giúp cùng một chính sách đọc được đặc trưng có ý nghĩa so sánh.

\subsection{Bộ trích đặc trưng}
% ~450 từ -- so sánh 4 phương án (ĐÓNG GÓP 3), mỗi phương án 1 đoạn ngắn:
% (a) MLP: làm phẳng ma trận nút, không dùng cấu trúc đồ thị (mốc dưới).
% (b) MLP + đặc trưng đồ thị: nối thêm trung bình đặc trưng lân cận 1-hop và
%     thống kê mức đồ thị -> có thông tin cấu trúc nhưng không cần lan truyền.
% (c) GCN: 3 lớp tích chập (32->64->128->256) trên ma trận kề đối xứng đã chuẩn
%     hoá; trọng số cạnh w_uv = exp(-d_uv) (đường ngắn -> liên kết mạnh); có kết
%     nối tắt; đọc ra theo nút rồi nối với nhánh MLP của đặc trưng toàn cục.
% (d) Attention: các khối tự chú ý đa đầu + gộp PMA (coi các nút là một tập).
% Fig.~1 = sơ đồ kiến trúc 2 nhánh (đồ thị + toàn cục) -> hợp nhất -> Q-head.

\subsection{Huấn luyện}
% ~250 từ + Bảng I siêu tham số. DQN (Stable-Baselines3): buffer 20k,
% batch 128, lr 1e-4 (tuỳ chọn giảm tuyến tính), epsilon 1.0 -> 0.05 trong 30%
% số bước, target update 1000, clip gradient 1.0, 200k bước, seed cố định,
% torch.use_deterministic_algorithms(True). Nêu phần cứng + thời gian huấn luyện.

\begin{algorithm}
\caption{Một episode mở rộng kế hoạch trạm sạc}
\label{alg:rollout}
\begin{algorithmic}[1]
\STATE Khởi tạo $P \leftarrow P_0$, ngân sách $B$, $\mathcal{S}^{*} \leftarrow \mathcal{S}(P_0)$
\WHILE{$B > 0$ và $t < |V|/2$}
  \STATE Quan sát $s_t$ (đặc trưng nút, cạnh, toàn cục)
  \STATE $a_t \leftarrow \arg\max_a Q_\theta(s_t,a)$ \COMMENT{$\varepsilon$-greedy khi huấn luyện}
  \STATE Chọn nút theo heuristic ứng với $a_t$; cập nhật $P$ và $B$
  \STATE Gán lại nút--trạm; cập nhật phụ tải thanh cái
  \STATE Tính $r_t$ theo (\ref{eq:reward}); cập nhật $\mathcal{S}^{*}$ và $P^{*}$
\ENDWHILE
\STATE \textbf{return} $P^{*}$
\end{algorithmic}
\end{algorithm}

%======================================================================
\section{Thiết lập thực nghiệm}
% NGÂN SÁCH: 0,7 trang.

\subsection{Dữ liệu}
% Bảng II: thống kê 4 quận (số nút, số cạnh, dân số, số trạm hiện hữu, số thanh
% cái 22 kV, ngân sách). Nguồn: OSM + dữ liệu TBA Hà Nội (500/220/110 kV).

\subsection{Các phương pháp đối sánh}
% 4 mốc so sánh, mỗi mốc 2--3 câu:
% - Kế hoạch hiện hữu (không can thiệp) -- mốc gốc.
% - Greedy-Benefit: mỗi bước chọn nút có lợi ích biên trên chi phí cao nhất.
% - Greedy-Demand: mỗi bước chọn nút có nhu cầu chưa đáp ứng lớn nhất.
% - GA: tiến hoá TRỌNG SỐ của mạng chính sách (neuroevolution) trên cùng không
%   gian quan sát MLP; fitness = điểm chuẩn hoá tốt nhất trong episode; nêu
%   pop_size, số thế hệ, tỉ lệ đột biến/lai ghép.
% Nhấn mạnh: tất cả khởi động từ CÙNG kế hoạch hiện hữu, cùng ngân sách,
% cùng đồ thị, cùng seed.

\subsection{Độ đo}
% Liệt kê chính xác các cột sẽ báo cáo:
% điểm chuẩn hoá; lợi ích; chi phí; công bằng; thời gian chờ LỚN NHẤT (phút);
% thời gian đi lại LỚN NHẤT (phút); phạt lưới; SỐ THANH CÁI QUÁ TẢI; số trạm;
% tỉ lệ ngân sách đã dùng.
% CẢNH BÁO: code hiện tính max chứ không phải trung bình -> phải ghi rõ "lớn
% nhất/xấu nhất", đừng viết "trung bình".

%======================================================================
\section{Kết quả và thảo luận}
% NGÂN SÁCH: 1,0--1,1 trang gồm Bảng III + Fig.~2 + Fig.~3.

% Bảng III (chính): so sánh 4 phương pháp x 4 quận. Nếu chật, chỉ để 2 quận
% trong bảng chính, 2 quận còn lại tóm tắt bằng 1 câu.
% -> BẮT BUỘC: báo cáo trung bình +/- độ lệch chuẩn trên >= 5 seed. Hiện repo
%    mới chạy 1 seed; đây là việc cần làm trước khi nộp.

% Fig.~2: bản đồ quận với vị trí trạm do RL chọn (Results/maps/*), đặt cạnh
%         kết quả của greedy để thấy khác biệt về phân bố không gian.
% Fig.~3: đường cong huấn luyện (reward + best score) HOẶC biên Pareto
%         (plots/pareto_frontier) -- chọn 1, không đủ chỗ cho cả hai.

\subsection{So sánh với các phương pháp đối sánh}
% 2 đoạn: (i) RL thắng ở đâu, thua ở đâu; (ii) GIẢI THÍCH VÌ SAO -- greedy tối
% ưu lợi ích tức thời nên dồn trạm vào cùng vùng phụ tải và làm quá tải thanh
% cái; chính sách RL học được cách "trả giá" ngắn hạn để giữ dung lượng lưới.

\subsection{Phân tích loại bỏ (ablation)}
% Bảng IV hoặc gộp vào Bảng III. 4 kịch bản đã có sẵn số liệu:
% Full / beta=0 (bỏ suy giảm theo khoảng cách) / eta=0 (bỏ co giãn nhu cầu) /
% bỏ phạt lưới / nhu cầu tĩnh.
% Thông điệp chính (số liệu hiện có ủng hộ): bỏ phạt lưới hoặc bỏ suy giảm nhu
% cầu đều làm TĂNG điểm thô nhưng kéo theo 1--2 thanh cái quá tải và 49 trạm
% (dùng gần 100% ngân sách) -- tức là mô hình đầy đủ đánh đổi điểm lấy tính khả
% thi vận hành. Phải diễn giải theo hướng này, đừng né việc điểm thô thấp hơn.

\subsection{Thảo luận về biểu diễn trạng thái}
% 1 đoạn ngắn so sánh MLP / MLP+graph / GCN / attention: cái nào đáng công sức.

%======================================================================
\section{Kết luận}
% NGÂN SÁCH: 0,2 trang (~150 từ). 3 câu tổng kết + 2 câu hạn chế (lưới điện
% tổng hợp từ toạ độ TBA chứ chưa có số liệu vận hành thực; một thành phố; ngân
% sách cố định) + 2 câu hướng phát triển (đa mục tiêu, dữ liệu lưới thực,
% chuyển giao chính sách giữa các quận/thành phố).

%======================================================================
\section*{Cảm tạ}
% 1--2 câu (đề tài/quỹ tài trợ). Cắt bỏ nếu thiếu trang.

%======================================================================
% TÀI LIỆU THAM KHẢO -- NGÂN SÁCH: 0,4 trang, 15--18 mục.
% Phân bổ gợi ý: 5 định vị trạm sạc, 3 GA/metaheuristic, 4 DRL + DRL trên đồ
% thị, 3 ràng buộc lưới/hosting capacity, 2 công cụ (Stable-Baselines3, OSMnx,
% pandapower).
\begin{thebibliography}{1}
\bibitem{sb3} A.~Raffin \emph{et al.}, ``Stable-Baselines3: Reliable reinforcement learning implementations,'' \emph{JMLR}, 2021.
\end{thebibliography}

\end{document}
